\section{Experiments and Results}
\label{sec:experiments}

\subsection{Experimental Setup}
To evaluate the efficacy and robustness of the proposed PC-Diffusion framework, experiments are conducted utilizing two widely recognized benchmark datasets: the Case Western Reserve University (CWRU) bearing dataset \cite{cwru2023bearing} and the Xi'an Jiaotong University - Sumyoung (XJTU-SY) bearing fault lifecycle dataset \cite{wang2018xjtu}. 

The CWRU dataset provides vibration signals collected at a sampling frequency of 12kHz under varying motor loads. The dataset encompasses 10 fault categories, including Normal and 3 fault locations (inner race, outer race, ball) with 3 severities (0.007, 0.014, and 0.021 inches). Each class contains approximately 500 samples, with a fixed sequence length of 1024 points. In contrast, the XJTU-SY dataset reflects accelerated run-to-failure conditions sampled at 25.6kHz, representing extreme industrial degradation over complex life-cycles under 5 running conditions. For both datasets, the samples are explicitly divided into training, validation, and testing sets with a ratio of 70\%, 10\%, and 20\%, respectively.

\subsection{Dataset Configuration for Zero-Shot Learning}
For both datasets, we formulate a Generalized Zero-Shot Learning (GZSL) protocol by splitting the fault conditions into explicitly disjoint seen and unseen categories. For the CWRU dataset, the seen classes consist of Normal, IR007, OR007, and Ball007 (4 classes), while the unseen classes are IR014, OR014, Ball014, IR021, OR021, and Ball021 (6 classes). For the XJTU-SY dataset, the split is performed based on running conditions, where conditions 1-3 are defined as seen classes, and conditions 4-5 are defined as unseen classes. During training, the PC-Diffusion model and the Bayesian embedding module only access the raw time-series signatures and semantics of the seen classes. The detailed split configuration is summarized in Table \ref{tab:zsl_split}.

\begin{table}[ht]
\centering
\caption{Zero-Shot Learning Split Configuration for Experimental Datasets}
\label{tab:zsl_split}
\begin{tabular}{lll}
\toprule
\textbf{Dataset} & \textbf{Seen Classes (Training/Testing)} & \textbf{Unseen Classes (Testing Only)} \\
\midrule
CWRU & Normal, IR007, OR007, Ball007 (4 classes) & IR014, OR014, Ball014, IR021, OR021, Ball021 (6 classes) \\
XJTU-SY & Running Conditions 1-3 & Running Conditions 4-5 \\
\bottomrule
\end{tabular}
\end{table}

\subsection{Implementation Details}
Data preprocessing comprises standardized min-max normalization alongside a fixed temporal slicing window to yield uniformly sized input sequences of 1024 points. The PC-Diffusion architecture utilizes an enhanced WaveNet-based 1D convolutional network as the primary noise prediction back-bone. The models are trained with a batch size of 64 using the Adam optimizer with a learning rate of $2 \times 10^{-4}$ and a weight decay of $1 \times 10^{-5}$. The total training comprises 200 epochs, which includes 50 epochs for PINN pre-training and 150 epochs for training the PC-Diffusion model. For semantic generation, GPT-4 is employed as the underlying LLM to generate comprehensive and accurate fault semantic descriptions for the Bayesian embedding. The entire infrastructure is implemented in PyTorch 2.1 and accelerated utilizing dual RTX 4090 GPUs.

\subsection{Quality of Generated Samples}
A primary comparative metric concerns the physical viability of the zero-shot generated unseen samples. We compare synthetic vibrations produced by a canonical DDPM \cite{ho2020denoising} against our PC-Diffusion framework. Fast Fourier Transform (FFT) analysis dictates that conventional diffusion models generate erratic spectral spikes entirely disassociated with the theoretical Outer Race Fault Frequency (BPFO) or Inner Race Fault Frequency (BPFI). Conversely, samples drawn from PC-Diffusion exhibit remarkable spectral energy alignment with theoretical fault kinematics, verifying that $\mathcal{L}_{physics}$ successfully bounds the generative vector space to reality.

\subsection{Zero-Shot Diagnosis Results}
Under the GZSL protocol, rigorous metrics including Seen Accuracy ($Acc_S$), Unseen Accuracy ($Acc_U$), and their Harmonic Mean ($H$) are reported in Table \ref{tab:performance}. Compared to classical semantic autoencoder methods like CADA-VAE \cite{schonfeld2019cada} and VAE/GAN-driven GZSL models such as standard VAE \cite{kingma2013auto} and WGAN-GP \cite{gulrajani2017improved}, PC-Diffusion achieves an outstanding progression. On the XJTU-SY dataset, our model successfully records an $H$ score approximating 91.4\%, drastically out-performing baseline standards which oscillate near 75\%. The Bayesian mapping effectively prevents the characteristic degradation of seen-class diagnostic capability when the classifier is exposed to unseen categories. The detailed class-wise classification performance is visualized in the confusion matrix in Fig. \ref{fig:confusion}.

\subsection{Uncertainty Quantification Analysis}
To address measurement safety, we compute the predictive entropy across all diagnostic classifications. For instances derived from overlapping noise profiles or anomalous domain distributions, the Bayesian semantic module successfully quantifies an elevated variance threshold \cite{gal2016dropout}. By setting a predictive entropy threshold (e.g., $H>0.5$ is marked as uncertain), the framework avoids over-confident incorrect predictions regarding unseen data. Notably, the model achieves a seen class average confidence of 95.2\%$\pm$2.1\% and an unseen class average confidence of 87.6\%$\pm$4.3\%, returning highly distinguishable confidence bounds as shown in Fig. \ref{fig:uncertainty}. Furthermore, for out-of-distribution (OOD) detection, the framework achieves an impressive AUROC of 0.943. This capability ensures critical out-of-distribution faults are flagged for manual expert inspection rather than blindly misdiagnosed. 

\subsection{Ablation Studies}
Ablation configurations isolating the physically-guided diffusion process and the dual-level NLP semantic embeddings prove their necessity. Compared to the full model's Harmonic Mean score of 91.4\%, removing the frequency energy loss ($\mathcal{L}_{freq}$) prompts severe degradation to $H=79.1\%$ ($\Delta H = -12.3\%$). Similarly, omitting the time-domain loss ($\mathcal{L}_{time}$) results in a drop to $H=83.7\%$ ($\Delta H = -7.7\%$). Removing the Bayesian mapping yields $H=85.2\%$ ($\Delta H = -6.2\%$). Additionally, substituting the dual-level semantics with single-level semantics (without LLM) collapses the flexibility observed during unseen inference, reducing the score to $H=82.8\%$ ($\Delta H = -8.6\%$).

\subsection{Interpretability Analysis}
Deploying t-Distributed Stochastic Neighbor Embedding (t-SNE) \cite{vandermaaten2008tsne} validates that generated unseen features systematically cluster within anatomically correct spatial regions dictated by their analog semantic attributes, as illustrated in Fig. \ref{fig:tsne}. The boundaries between disparate synthetic fault states remain exceptionally tight and discernible.